\documentclass[11pt]{article}
\usepackage[utf8]{inputenc}
\usepackage[T1]{fontenc}
\usepackage[a4paper,margin=1in]{geometry}
\usepackage{setspace}
\usepackage{parskip}
\usepackage{hyperref}
\hypersetup{hidelinks}
\usepackage{authblk}

\title{\textbf{A Universe from Coulomb and Newton}\\[0.3em]
\large Two Mirror Forces --- Charge for the Small, Mass for the Large}
\author[ ]{Claes Johnson}
\affil[ ]{KTH Royal Institute of Technology, Stockholm \quad\texttt{claesjohnson@gmail.com}}
\date{\today}

\begin{document}
\maketitle
\setstretch{1.06}

\begin{abstract}
\noindent A non-technical companion to the article \emph{Coulomb $+$ Newton: Towards a Cosmology}. It asks whether the
universe can be built from just two forces that obey the \emph{same} law in mirror image --- Coulomb's law for
electric charge and Newton's law for gravity --- with no Big Bang singularity, no inflation and no dark-matter
particle. Coulomb, acting on charge, builds matter at the small scale; Newton, acting
on the energy of that matter, gathers it at the large scale. The price, stated plainly at the end, is that it
explains neither the accelerating expansion nor the spinning of galaxies.
\end{abstract}

\section*{All cosmology is speculative}
It is worth saying plainly at the start: every cosmology is speculative at its foundations. The early universe is
not something we can visit or repeat. The standard account --- a hot Big Bang, a burst of inflation, a cosmos
dominated by dark matter and dark energy --- fits the data, but it does so with ingredients that are inferred, not
seen: a singularity no instrument reaches, an inflationary epoch invented to smooth things out, and two ``dark''
substances that make up most of the universe yet have never appeared in any laboratory. In a subject resting on
such foundations, it is fair --- and useful --- to ask what a \emph{simple} alternative would look like, built not
from new and unseen things but from the most ordinary ingredients physics already has.

This essay sketches one such alternative. Its ingredients are two laws you already know --- \textbf{Coulomb's law}
for electric charge and \textbf{Newton's law} for gravity.

\section*{Two forces, mirror images of each other}
Here is the first thing worth noticing, and it is old news that we rarely dwell on: Coulomb's law and Newton's law
have \emph{the same shape}. Both fall off as one over distance squared; both can be written as the same field
equation. They differ in only two ways:

\begin{itemize}
\item \textbf{a sign.} For electric charge, \emph{like repels like} (two positives push apart). For gravity,
\emph{like attracts like} (two masses pull together). Gravity is Coulomb with the sign flipped.
\item \textbf{a strength.} Each force has its own constant fixing how strong it is --- and gravity's is
fantastically smaller.
\end{itemize}

So the two forces are mirror images: the same law, one attractive and one repulsive, each with its own strength.
This essay's proposal is simply that these two mirror forces, acting on two things --- \emph{charge} and
\emph{mass} --- are enough to build a universe.

\section*{Two sources, two scales}
The division of labour is clean. \textbf{Coulomb, acting on charge, rules the small}: it builds nuclei, atoms,
molecules --- all of matter and chemistry. \textbf{Newton, acting on mass, rules the large}: it gathers matter
into stars, galaxies and the cosmic web. Charge is the source of the one force; mass (which, as we will see, is
just energy) is the source of the other. Two forces, two sources, two scales --- running in parallel, neither
derived from the other.

\section*{What replaces the bang}
Every cosmology needs a beginning for its \emph{structure}, if not for time itself. The standard one is the Big
Bang: a hot dense state, with a brief episode of inflation stretching microscopic ripples into the seeds of
galaxies. This picture asks for something quieter. Suppose the electric potential --- the field whose slopes push
charges around --- is not perfectly flat but gently \textbf{rippled}, at every size of ripple at once.

Then two things happen at once, at opposite ends of the scale. Where the ripples are finest, the field curves
sharply, and it is curvature that makes charge. Each tiny crest becomes a unit of positive charge, each trough a
unit of negative: protons and electrons, in exactly equal numbers, because every crest has its trough. The
graininess of matter is the graininess of the ripples.

Here something pleasing happens to the arithmetic. One might expect the \emph{size} of the ripples to be a free
knob, tunable to taste. It isn't. Demanding that each ripple carry exactly \emph{one} charge --- no halves, no
doubles --- fixes how strong it must be, once you have said how small it is. And how small it is turns out to be
the same unanswered question as how big a proton is. So the entire small-scale story runs on a single unknown
length, with everything else following.

The coarsest ripples do something quite different. They cannot be felt as electric forces --- a very gradual slope
in the potential is invisible to matter, and matter rearranges itself to cancel any large-scale imbalance anyway.
Instead their job is subtler: to govern \emph{where the fine ripples are busiest}, so that matter ends up more
abundant here and sparser there. That unevenness is what gravity later works on, and it is where the cosmic web
comes from.

That last step is the one genuine gamble in the picture, and it is worth flagging rather than gliding past. It
only works if the rippling is \emph{lopsided} --- if the broad swells and the fine ripples are correlated rather
than independent. A perfectly even-handed, statistically featureless rippling would leave matter uniformly spread,
and there would be no structure at all. Curiously, the same lopsidedness is needed elsewhere in the picture, to
explain why the positive charges come out small and dense and the negative ones large and diffuse. One assumption,
two debts paid. Whether it is true is a calculation nobody has yet done.

\section*{The small world: Coulomb builds matter}
Electric charge comes in two signs, and in this picture they come out \emph{different in size}. The positive
charge forms small, dense clumps --- \textbf{protons}. The negative charge spreads into large, diffuse clouds ---
\textbf{electrons}. From just these, Coulomb's law builds the rest by simple geometry: electrons bound around
protons make \textbf{atoms}; electrons shared between atoms make \textbf{molecules} and all of chemistry.

Remarkably, the same force reaches inside the nucleus. In this picture a nucleus is not held together by a separate
``strong force.'' It is protons glued by \emph{compact} electrons --- the very same electrons, squeezed small ---
again by ordinary Coulomb attraction, only at a tiny scale where it is correspondingly strong. What looks like a
new and powerful nuclear force is just electricity at close range. One force, two scales.

\section*{The large world: Newton gathers matter}
Now the other mirror force. The matter Coulomb has built carries \emph{energy}, and energy is what gravitates (a
word on that below). Newton's law acts on it with the mirror of Coulomb's rule: where like charges \emph{repel},
like masses \emph{attract}.

But there is a subtlety here, and it is the heart of the picture. Gravity does not respond to how much matter
there is; it responds to how much matter there is \emph{compared with the average}. A region denser than average
pulls its surroundings in. A region emptier than average does the opposite: it \emph{pushes} them away. So even
though every scrap of matter is ordinary and positive, the thing gravity actually listens to comes in \textbf{two
signs} --- surplus and shortfall --- and the mirror rule then does its work. Surplus gathers with surplus, and the
matter drains out of the shortfalls.

That is the whole story of the large scale. The surpluses become the glowing \textbf{cosmic web} of filaments and
clusters that we live in; the shortfalls become the great \textbf{voids} between them --- and the voids are not
merely leftover empty space, but actively push their surroundings outward and so grow. No negative matter is
needed anywhere: nothing falls upward, no strange substance is added. Emptiness relative to the average is
repulsion enough.

One thing this picture does \emph{not} give, and it is only honest to say so at once: it does not explain the
accelerating expansion that astronomers attribute to dark energy. Voids push locally, which is real, but it is a
local effect, and ordinary cosmology has it too. Across the universe as a whole the surpluses and the shortfalls
cancel exactly --- that is what ``average'' means --- so the expansion neither slows nor speeds up. It
\textbf{coasts}. That is a definite claim, and a risky one, because it can be checked against exploding stars
across the sky; we return to it at the end.

\section*{The web as unmixing}
There is a familiar name for what this sorting really is. Shake oil and water together, let the jar stand, and the
two do not stay mixed: they \emph{separate}, oil gathering with oil and water with water, until thin skins divide
them into a pattern of blobs and films. Physicists call this \textbf{phase separation}, and it is one of the
best-understood pattern-forming processes in nature --- the very same thing that happens when a hot alloy cools and
its two metals unmix into a marbled web.

The universe of surplus and shortfall does exactly this. Surplus gathers with surplus, and matter drains from the
shortfalls, until the two are separated into a marbled pattern of filaments threaded through great empty rooms.
Seen this way the glowing cosmic web is not a special or mysterious shape at all: it is the ordinary shape that
things take when each gathers with its own kind --- the same web you would get from unmixing oil and water, only
drawn by gravity across the whole sky.

The analogy has one limit worth marking, because it is easy to overdraw. Oil and water are symmetric: swap them
and the picture looks the same. Surplus and shortfall are \emph{not}. A filament can go on getting denser without
limit, but a void can at most become completely empty --- there is a floor to emptiness and no ceiling to
crowding. So the two are not mirror images of one another, and one should not expect voids to be shaped like
filaments turned inside out. What halts the clumping before everything collapses to points is the ordinary
springiness of matter: squeeze a gas and it heats, and past a certain stiffness the heating wins.

\section*{A word about ``mass''}
We have been saying gravity acts on \emph{mass}, then on \emph{energy}, as if they were the same. They are. Mass,
looked at closely, is just a name for \emph{energy content}: a thing's resistance to being pushed is set by how
much energy it holds --- that is the real meaning of $E=mc^2$, and it is measured directly (a single particle can
be weighed by its energy; a nucleus weighs less than its parts by exactly the energy of their binding). So gravity
acts on the energy of matter, which is what we usually call its mass. We need no mysterious extra substance called
mass; energy suffices.

One further temptation should be resisted, honestly. It is tempting to say a particle's energy is fixed
\emph{entirely} by the size of its charge cloud --- small and tight meaning high energy, like a stiff spring.
Qualitatively that is surely right (the small proton is the heavy one). But the exact numbers do not line up: the
proton is more compact than a simple size--energy rule would make it, by about a factor of ten. So we say only what
is safe --- energy and size are closely \emph{related}, not exactly \emph{identified} --- and we do \emph{not}
claim that gravity is somehow secretly derived from electricity. The two forces run side by side; whether they are
ultimately one is an open question, not a claim of this picture.

\section*{Three numbers}
Strip the picture to essentials and it rests on remarkably little: \emph{three pure numbers}, each doing one job.
The first is the strength of electricity --- what physicists call the fine-structure constant, about $1/137$ ---
and it fixes how much bigger an atom is than its nucleus: about twenty thousand times. The second is how much
heavier a proton is than an electron: $1836$. Those two, between them, settle all of atoms, chemistry and nuclei,
and neither of them mentions gravity at all.

The third is the strength of gravity, and it is the strangest number in physics: measured against electricity, it
is smaller by roughly $10^{39}$. Pile up matter and gravity is negligible --- until you have gathered something
like $10^{57}$ protons, which is to say about the mass of the Sun, and then it wins. That single number is what
separates the small world from the large one. And its absurd smallness is not an embarrassment: it is exactly why
there is room for chemistry, for planets, for everything between an atom and a star. Were gravity anywhere near as
strong as electricity, there would be no in-between at all, and nothing here to describe.

Compare the bookkeeping. The standard model of cosmology needs a singularity, an inflaton field, a dark-matter
particle and a cosmological constant --- a lengthening list of unseen ingredients. This picture asks for three
numbers. One qualification, since it is easy to overstate: a term of the same kind as the cosmological constant
does appear here, in the rule that only departures from the average gravitate. The difference is that it is not
free to be chosen --- it is fixed by the requirement that the universe's gravity cancels overall, and it thins out
as the universe expands instead of staying fixed. The open questions are correspondingly few and sharp: none of the three is derived from anything deeper
(nor are they in the standard account), and whether the two forces are at bottom one remains genuinely open. Three
clean unknowns, in place of a menagerie of invisible substances --- though, as the next section admits, the price
is paid elsewhere.

\section*{A comparison with the standard model, and a soft spot in it}
The standard cosmology --- $\Lambda$CDM --- rests on a hot Big Bang, with the cosmic microwave background (CMB) as
its relic glow, supplemented by dark matter, dark energy and an early burst of inflation. Its single proudest piece
of evidence is the CMB's spectrum: an almost perfect \emph{Planck} (blackbody) curve, read as the thermal radiation
of a hot, dense early universe once in equilibrium.

But this reads the evidence backwards. A perfect Planck spectrum is not something a hot dense gas naturally makes.
In the author's analysis of blackbody radiation, a true Planck curve is the signature of a \emph{resonant lattice}
--- a regular array of oscillators (a metallic crystal, a \emph{gitter}) with a sharp high-frequency cutoff --- and
\emph{not} of the disordered collisions of a hot gas or plasma, which have no such lattice and no such cutoff. On
that analysis a hot dense fireball simply \emph{cannot} generate a perfect Planck spectrum. (This contests the
textbook view, by which any body in thermal equilibrium is supposed to radiate Planckian whatever its make-up; it
is the author's position, developed at length in the blackbody work cited below.)

If so, the conclusion is sharp: \textbf{the CMB's near-perfect Planck spectrum is exactly what a hot dense gas would
\emph{not} produce} --- so it cannot serve as evidence that any such hot dense initial state ever existed.
\textbf{The CMB can say nothing about the existence of a hot dense beginning.} Its Planck form points, if anything,
toward a resonant, lattice-like source --- whatever and wherever that turns out to be --- not toward a primordial
fireball.

This is where the present picture simply steps aside. It has \emph{no} hot Big Bang and no primordial fireball, so
it neither needs nor claims the CMB as the relic of one. What the CMB is, and why it is so smooth and so Planckian,
is left here as a separate open question --- but its spectrum is \emph{not} read as proof of a hot dense
beginning, because on the blackbody analysis above it is no such proof.

\section*{A universe you could almost draw}
That is the appeal. In place of an unrepeatable explosion followed by a cascade of unseen ingredients, here is a
universe you could nearly sketch on a napkin: two mirror forces, one repulsive and one attractive, obeying the same
law; electric charge of two sizes --- small dense protons, large diffuse electrons; Coulomb assembling them into
nuclei, atoms and chemistry; and Newton, acting on the energy of all that matter --- on its surpluses and
shortfalls, not on its bulk --- gathering it into a cosmic web of filaments and clusters threaded by voids that
push outward as they empty. Two old laws, three numbers.

It should be said just as plainly what this costs. The picture does not explain the accelerating expansion that
dark energy was invented for --- it says the universe coasts instead, which is a different claim and a checkable
one. It does not explain why galaxies spin as fast at their edges as at their centres, the observation that dark
matter was invented for; that problem is left exactly where Newton left it. And it does not yet explain why a
nucleus is as small as it is, missing by a factor of about thirty. Three debts, stated openly.

So the simplest question of all is left standing: is nature really this economical? It may not be. But it is worth
knowing what the simple answer looks like --- and precisely where it fails --- before reaching for the complicated
one.

\bigskip
\hrule
\medskip
\noindent\emph{For the mathematics --- the two-valued Poisson equations, the RealQM / RealNucleus computations, the
cosmic web of overdense filaments and underdense voids, and the quantitative claims and open problems only gestured
at here --- see the companion technical article,} C.\ Johnson, \emph{Coulomb $+$ Newton: Towards a Cosmology},
\emph{and the treatise} \emph{Real Quantum Mechanics} \emph{($\sim$340 pp), both free from
the RealQM gallery,} \url{https://claes542.github.io/RealMolecule/}.

\medskip
\noindent\emph{For the blackbody argument --- that a perfect Planck spectrum is the signature of a resonant lattice,
not of a hot dense gas --- see} C.\ Johnson, \emph{Computational Blackbody Radiation} \emph{(online book).}

\end{document}
